\begin{Exercise}[label={ex:ms}]
    Utilizzando 8 bit, rappresentare il numero decimale $-105_{10}$ in Modulo e Segno (MS).
\end{Exercise}

\begin{Answer}[ref={ex:ms}]
    Il numero è negativo, quindi il bit di segno sarà 1.
    Calcoliamo con i restanti 7 bit il valore assoluto, $|-105| = 105_{10}$:
    
    \begin{center}
    \begin{tabular}{r|r}
        \multicolumn{2}{c}{2} \\
        105 & 1 \\
        52 & 0 \\
        26 & 0 \\
        13 & 1 \\
        6 & 0 \\
        3 & 1 \\
        1 & 1 \\
        0 &
    \end{tabular}
    \end{center}

    $105_{10} = 1101001_2$
    
    Unendo il bit di segno al valore assoluto convertito in base due ottienamo come risultato finale $-105_{10} = 11101001_{2\text{MS}}$
\end{Answer}
